% Original analytic teaching diagram, 2026-11-09.
\definecolor{VB}{HTML}{7188A3}\definecolor{VBF}{HTML}{E5EAF0}
\definecolor{VG}{HTML}{829C8B}\definecolor{VGF}{HTML}{E8EEE9}
\definecolor{VY}{HTML}{B4A16F}\definecolor{VYF}{HTML}{F1ECDD}
\definecolor{VR}{HTML}{A7787D}\definecolor{VRF}{HTML}{F0E5E6}
\definecolor{VP}{HTML}{9588A5}\definecolor{VPF}{HTML}{ECE8F0}
\definecolor{VI}{HTML}{3B4252}\definecolor{VM}{HTML}{6B7280}\definecolor{VL}{HTML}{4C4D4F}
\begin{tikzpicture}[x=.95mm,y=1mm,>=Stealth,font=\figurefont\fontsize{8}{10}\selectfont,
lab/.style={text=VM,font=\figurefont\fontsize{7.5}{9}\selectfont},
box/.style={draw=VB,fill=VBF,line width=.7pt,rounded corners=1mm,minimum height=8.5mm,minimum width=34mm,align=center,text=VI},
stop/.style={draw=VR,fill=VRF,line width=.7pt,rounded corners=1mm,minimum height=8.5mm,minimum width=38mm,align=center,text=VI},
ok/.style={draw=VG,fill=VGF,line width=.7pt,rounded corners=1mm,minimum height=8.5mm,minimum width=38mm,align=center,text=VI},
ar/.style={->,draw=VL,line width=.75pt}]
\node[box] (c) at (0,0) {方向$\bm{d}^{(j)}$的曲率\\$\kappa=\bm{d}^{(j)\top}\bm{B}\bm{d}^{(j)}$};
\node[box] (b) at (0,-17) {试走到\\$\bm{p}^{(j)}+\alpha^{(j)}\bm{d}^{(j)}$};
\node[box] (r) at (0,-34) {更新残差\\$\bm{r}^{(j+1)}$};
\node[box] (n) at (0,-51) {新方向$\bm{d}^{(j+1)}$};
\node[stop] (s1) at (58,0) {曲率停止\\沿$\bm{d}^{(j)}$走到边界};
\node[stop] (s2) at (58,-17) {边界停止\\截在边界上};
\node[ok] (s3) at (58,-34) {残差停止\\返回当前$\bm{p}^{(j+1)}$};
\draw[ar] (c)--node[left,lab]{$\kappa>0$} (b);
\draw[ar] (b)--node[left,lab]{仍在区域内} (r);
\draw[ar] (r)--node[left,lab]{残差仍大} (n);
\draw[ar] (c)--node[above,lab]{$\kappa\leq0$} (s1);
\draw[ar] (b)--node[above,lab]{越出区域} (s2);
\draw[ar] (r)--node[above,lab]{残差足够小} (s3);
\draw[ar] (n.west)--++(-5,0)|-node[pos=.25,left,lab]{$j\leftarrow j+1$} (c.west);
\end{tikzpicture}
